\documentclass{article}
\begin{document}
\section{Introduction}
Citations to many sources \cite{s1,s2,s3,s4,s5,s6,s7,s8,s9,s10,s11,s12} break across lines
like any other text, and the reference list numbers them to twelve, so its labels are as wide as
the label of the widest, given as \texttt{99}. Later entries are cited alone \cite{s11} and with a
note \cite[remark]{s12}.
Pages are built one line at a time. As each line arrives the program measures how well the page would look if it ended there, and when adding more would only make it worse it stops and starts a new page. Penalties discourage breaks in awkward places, such as after the first line of a paragraph or before its last.

Vertical space between paragraphs and around headings can stretch and shrink a little, so that a page whose lines do not exactly fill it can still be set with its last line at the bottom. The amount of stretch is small, and most of it goes where a reader will not notice.
\begin{thebibliography}{99}
\bibitem{s1} A. Author1. A study of item number 1 in a long list of references, with enough words to need a second line. Journal of Examples 1, 1961.
\bibitem{s2} A. Author2. A study of item number 2 in a long list of references, with enough words to need a second line. Journal of Examples 2, 1962.
\bibitem{s3} A. Author3. A study of item number 3 in a long list of references, with enough words to need a second line. Journal of Examples 3, 1963.
\bibitem{s4} A. Author4. A study of item number 4 in a long list of references, with enough words to need a second line. Journal of Examples 4, 1964.
\bibitem{s5} A. Author5. A study of item number 5 in a long list of references, with enough words to need a second line. Journal of Examples 5, 1965.
\bibitem{s6} A. Author6. A study of item number 6 in a long list of references, with enough words to need a second line. Journal of Examples 6, 1966.
\bibitem{s7} A. Author7. A study of item number 7 in a long list of references, with enough words to need a second line. Journal of Examples 7, 1967.
\bibitem{s8} A. Author8. A study of item number 8 in a long list of references, with enough words to need a second line. Journal of Examples 8, 1968.
\bibitem{s9} A. Author9. A study of item number 9 in a long list of references, with enough words to need a second line. Journal of Examples 9, 1969.
\bibitem{s10} A. Author10. A study of item number 10 in a long list of references, with enough words to need a second line. Journal of Examples 10, 1970.
\bibitem{s11} A. Author11. A study of item number 11 in a long list of references, with enough words to need a second line. Journal of Examples 11, 1971.
\bibitem{s12} A. Author12. A study of item number 12 in a long list of references, with enough words to need a second line. Journal of Examples 12, 1972.
\end{thebibliography}
\end{document}
